\documentclass[11pt, a4paper]{article}

%\usepackage{geometry}
%\usepackage{epsfig}
%\usepackage{epstopdf}
%\usepackage{graphics}
\usepackage{booktabs}
\usepackage[svgnames]{xcolor}
\usepackage{graphicx}
\usepackage{hyperref}
\usepackage{parskip}
\usepackage{url}

\setlength{\oddsidemargin}{0cm}
\setlength{\evensidemargin}{0cm}
\setlength{\topmargin}{-1cm}
\setlength{\textheight}{23cm}
\setlength{\textwidth}{16cm}

%\setlength{\parindent}{0pt}

\pagestyle{headings}

\newcommand{\bat}{{\sc BAT}}
\newcommand{\BAT}{\bat}
\newcommand{\Root}{{\sc Root}}

\newcommand{\versionno}{1.1.0-DEV}

\newcommand{\version}{version~\versionno}
\newcommand{\Version}{Version~\versionno}

\newcommand{\code}[1]{\texttt{#1}}

\newcommand{\dan}[1]{{\color{red}{Dan: #1}}}
\newcommand{\fred}[1]{{\color{blue}{Fred: #1}}}

%--------------------------------------------------------

\begin{document}

% --------------------------------------------------------
% title
% --------------------------------------------------------

\thispagestyle{empty}

\begin{figure}
\includegraphics[scale=0.25]{bat}
\end{figure}

\vspace*{1cm}

\begin{center}


{\Large An introduction to \bat}
\\

\vspace{1cm}

{\large \bat\ \version}

\end{center}

\thispagestyle{empty}

\vfill

\begin{center}
\today
\end{center}

\pagebreak

% --------------------------------------------------------
% table of contents
% --------------------------------------------------------

\thispagestyle{empty}

\enlargethispage{2cm}

\tableofcontents

\section{Intro to Bayesian Statistics (for our purposes)}
\fred

\section{Intro to Markov Chain Monte Carlo (for our purposes)}
\fred

R value.

\section{Basic BAT}

Goal of this section is that if you can read and test as you go along, and end up with functioning code.

\subsection{Non-code parts of BAT}
\dan

bat-project, bat-config

\subsection{Absolutely necessary requirements of your model}
\dan

Parameters, LogLikelihood, some form of setting a prior.

\subsection{Running the analysis}
\dan

Functions to call. How to output. Basic tuning. How to interpret runtime log output \& trouble shooting (nonconvergence, poor efficiencies). How to interpret summary log output.

\subsection{Plots}
\dan

Interpreting the basic output: marginals; parameters; knowledge update; correlation matrix.

Options for the different output types.


\section{Non-MCMC parts of BAT}

\subsubsection{Integration}
\fred

\subsubsection{Optimization}
\fred

\section{Accompanying Models}

\subsection{Using BAT's native data structures}

\subsection{base}

\subsection{mtf}


\section{Advanced BAT}

\subsection{ModelManager for Model Comparison}
\fred

\subsection{Defining your own factorized prior}
\dan

\subsection{Sharing / Loading samples}
\dan

\subsection{Performance / Threading / User's attention to thread-safety}

tips.

\section{Examples}

\section{Validation}

Proof we've checked BAT works.

\end{document}
